%%%%%%%% ICML 2021 EXAMPLE LATEX SUBMISSION FILE %%%%%%%%%%%%%%%%%

\documentclass{article}

% Recommended, but optional, packages for figures and better typesetting:
\usepackage{microtype}
\usepackage{graphicx}
\usepackage{subfigure}
\usepackage{booktabs} % for professional tables

% hyperref makes hyperlinks in the resulting PDF.
% If your build breaks (sometimes temporarily if a hyperlink spans a page)
% please comment out the following usepackage line and replace
% \usepackage{icml2021} with \usepackage[nohyperref]{icml2021} above.
\usepackage{hyperref}

% Attempt to make hyperref and algorithmic work together better:
\newcommand{\theHalgorithm}{\arabic{algorithm}}

% Use the following line for the initial blind version submitted for review:
\usepackage[accepted]{comp547}

% If accepted, instead use the following line for the camera-ready submission:
%\usepackage[accepted]{icml2021}

% The \icmltitle you define below is probably too long as a header.
% Therefore, a short form for the running title is supplied here:
\icmltitlerunning{Project Proposal}

\begin{document}

\twocolumn[
\icmltitle{Project Title}

% It is OKAY to include author information, even for blind
% submissions: the style file will automatically remove it for you
% unless you've provided the [accepted] option to the icml2021
% package.

% List of affiliations: The first argument should be a (short)
% identifier you will use later to specify author affiliations
% Academic affiliations should list Department, University, City, Region, Country
% Industry affiliations should list Company, City, Region, Country

% You can specify symbols, otherwise they are numbered in order.
% Ideally, you should not use this facility. Affiliations will be numbered
% in order of appearance and this is the preferred way.
\icmlsetsymbol{equal}{*}

\begin{icmlauthorlist}
\icmlauthor{Name Surname}{equal,to}
%\icmlauthor{Name Surname2}{equal,goo}
\end{icmlauthorlist}

\icmlaffiliation{to}{Department of Computer Engineering}
%\icmlaffiliation{goo}{Department of Mechanical Engineering}

\icmlcorrespondingauthor{Name Surname}{email}

% You may provide any keywords that you
% find helpful for describing your paper; these are used to populate
% the "keywords" metadata in the PDF but will not be shown in the document
%\icmlkeywords{Machine Learning, ICML}

\vskip 0.3in
]

% this must go after the closing bracket ] following \twocolumn[ ...

% This command actually creates the footnote in the first column
% listing the affiliations and the copyright notice.
% The command takes one argument, which is text to display at the start of the footnote.
% The \icmlEqualContribution command is standard text for equal contribution.
% Remove it (just {}) if you do not need this facility.

%\printAffiliationsAndNotice{}  % leave blank if no need to mention equal contribution
\printAffiliationsAndNotice{\icmlEqualContribution} % otherwise use the standard text.

\begin{abstract}
  
\end{abstract}

\section{Introduction}
Introduce the task that you are going to investigate in your course project. State why you find your project topic interesting and what is difficult about it.

\section{Related Work}
Review previous work most relevant to your project topic. Discuss how you might improve upon these existing approaches.

\section{The Approach}
Give a brief outline of your approach. Describe the architecture you will use, whether you will extend an existing implementation, etc. Please note that you can change your approach later. 

\section{Experimental Evaluation}
Explain which dataset(s) you will use to train and test your model. Describe how you will evaluate the performance of your approach against those of competing methods.

\section{Work Plan}
Provide a rough timeline about the planned activities and their approximate deadlines. For example,

\begin{tabular}{|l|l|}
\hline
\textbf{Activity} & \textbf{Deadline} \\ 
\hline
Complete the literature search &	 MM/DD/YY \\
\hline
Reproduce results of a baseline approach &	MM/DD/YY \\ 
\hline
Prepare progress report &	MM/DD/YY \\
\hline
Make improvements X, Y, Z	& MM/DD/YY \\
\hline
Prepare final report and presentation & MM/DD/YY \\
\hline
\end{tabular}

\cite{Hinton:2006,Goodfellow:2014}

\bibliography{example_proposal}
\bibliographystyle{comp547}

\end{document}


% This document was modified from the file originally made available by
% Pat Langley and Andrea Danyluk for ICML-2K. This version was created
% by Iain Murray in 2018, and modified by Alexandre Bouchard in
% 2019 and 2021. Previous contributors include Dan Roy, Lise Getoor and Tobias
% Scheffer, which was slightly modified from the 2010 version by
% Thorsten Joachims & Johannes Fuernkranz, slightly modified from the
% 2009 version by Kiri Wagstaff and Sam Roweis's 2008 version, which is
% slightly modified from Prasad Tadepalli's 2007 version which is a
% lightly changed version of the previous year's version by Andrew
% Moore, which was in turn edited from those of Kristian Kersting and
% Codrina Lauth. Alex Smola contributed to the algorithmic style files.
